\chapter{FUNDAMENTAÇÃO TEÓRICA E TRABALHOS RELACIONADOS}

Este capítulo apresenta os fundamentos teóricos necessários à compreensão do problema abordado e ao desenvolvimento de um simulador espacial baseado em agentes para representar aspectos da transmissão da dengue. Inicialmente, são discutidos os principais conceitos epidemiológicos relacionados à doença e ao seu vetor. Em seguida, são apresentados conceitos de modelagem e simulação computacional, com ênfase na modelagem baseada em agentes e nos modelos epidemiológicos representados por estados.

Também são abordados a representação espacial por meio de Sistemas de Informação Geográfica (SIG), os conceitos de estocasticidade e experimentação computacional e os recursos da plataforma GAMA e da linguagem GAML. Por fim, são discutidos aspectos relacionados à documentação, verificação e validação de modelos, além de trabalhos anteriores relacionados à aplicação de simulações baseadas em agentes em problemas epidemiológicos.

\section{Fundamentos epidemiológicos da dengue}

A dengue é uma arbovirose causada pelo vírus da dengue e transmitida principalmente por mosquitos do gênero \textit{Aedes}. A compreensão dos mecanismos básicos de transmissão é necessária para definir quais entidades, estados e interações devem ser considerados em uma representação computacional do fenômeno.

\subsection{Doença, agente etiológico e sorotipos}

O vírus da dengue apresenta quatro sorotipos, denominados DENV-1, DENV-2, DENV-3 e DENV-4. A infecção pode ser assintomática ou apresentar manifestações clínicas de diferentes níveis de gravidade. A circulação de diferentes sorotipos constitui um aspecto relevante da dinâmica epidemiológica da doença, especialmente devido às relações entre infecções anteriores, resposta imunológica e risco de formas graves \cite{who_dengue}.

Do ponto de vista da modelagem computacional, a existência de múltiplos sorotipos demonstra que a representação da recuperação de um indivíduo pode envolver simplificações. Um modelo que represente apenas uma dinâmica epidemiológica genérica pode ser adequado ao seu propósito, desde que suas hipóteses sejam explicitadas e suas limitações sejam consideradas na interpretação dos resultados.

\subsection{Vetor e ciclo de desenvolvimento}

No contexto urbano brasileiro, o \textit{Aedes aegypti} possui papel central na transmissão da dengue. As fêmeas adultas participam diretamente da transmissão do vírus, enquanto o desenvolvimento do mosquito envolve estágios anteriores, como ovo, larva e pupa.

O ciclo de desenvolvimento do vetor é influenciado por condições ambientais, incluindo temperatura e disponibilidade de água. Dessa forma, a representação computacional do mosquito pode envolver tanto sua dinâmica epidemiológica quanto processos relacionados ao seu desenvolvimento e à disponibilidade de criadouros.

Entretanto, o nível de detalhamento adotado em um modelo deve ser compatível com seus objetivos, com os dados disponíveis e com o custo computacional da simulação. Nem todos os processos biológicos precisam ser representados individualmente para que o modelo possa investigar uma questão específica. A escolha das abstrações deve, portanto, ser explicitada durante a especificação do modelo.

\subsection{Transmissão humano--vetor}

A transmissão da dengue envolve interações entre humanos e mosquitos vetores. Um mosquito suscetível pode adquirir o vírus ao realizar alimentação sanguínea em um indivíduo infectante. Após um período de incubação extrínseca, o mosquito pode tornar-se capaz de transmitir o vírus. De maneira semelhante, um indivíduo suscetível pode adquirir a infecção após a picada de um mosquito infectante.

Esses processos evidenciam a necessidade de distinguir conceitos como infecção, período de incubação e infectividade. Um indivíduo ou vetor recém-infectado não necessariamente se torna imediatamente capaz de transmitir o vírus. Em modelos computacionais, essa característica pode ser representada por estados intermediários, permitindo distinguir agentes expostos daqueles que já se encontram em condição infectante.

A ocorrência da transmissão também depende das oportunidades de contato entre hospedeiros e vetores. Em modelos espacialmente explícitos, a localização, a mobilidade e as regras de proximidade podem influenciar a probabilidade de interação entre os agentes.

\subsection{Fatores ambientais e medidas de controle}

A dinâmica da dengue é influenciada por fatores biológicos, ambientais e sociais. Temperatura e precipitação podem afetar diferentes etapas do ciclo do vetor, enquanto a disponibilidade de criadouros e a distribuição espacial da população podem modificar as oportunidades de contato entre humanos e mosquitos.

Além disso, medidas de prevenção e controle podem atuar sobre diferentes componentes do sistema. A redução de criadouros, as intervenções destinadas ao controle vetorial e outras estratégias de prevenção não possuem necessariamente o mesmo mecanismo de ação. Dessa forma, quando representadas em um modelo computacional, essas medidas devem ser associadas a regras compatíveis com o efeito que se pretende investigar.

\section{Modelagem e simulação computacional}

A modelagem computacional consiste na construção de uma representação simplificada de um sistema com o objetivo de investigar determinados aspectos de seu funcionamento. A simulação corresponde à execução das regras definidas para o modelo ao longo do tempo, permitindo observar o comportamento produzido pelas relações entre seus componentes.

É importante distinguir o fenômeno real, o modelo conceitual, a implementação computacional e os resultados produzidos pela simulação. Diferenças entre os resultados esperados e os resultados observados podem estar relacionadas às hipóteses adotadas, aos dados utilizados, às regras implementadas ou aos procedimentos experimentais empregados.

\subsection{Abordagens de simulação}

Diferentes abordagens podem ser utilizadas para representar sistemas dinâmicos. A dinâmica de sistemas enfatiza relações agregadas entre variáveis, estoques e fluxos. A simulação a eventos discretos representa alterações de estado desencadeadas por eventos. Já a modelagem baseada em agentes representa entidades individuais e suas interações com outras entidades e com o ambiente \cite{macal_north_2010}.

\{aq deveria ter uma tabela\}

As abordagens não são necessariamente excludentes. Um modelo baseado em agentes pode utilizar estados epidemiológicos semelhantes aos empregados em modelos compartimentais e produzir indicadores agregados para análise.

\subsection{Modelos explicativos e preditivos}

A finalidade de um modelo influencia as abstrações utilizadas, os dados necessários e os critérios adotados para avaliar seus resultados. Modelos explicativos buscam investigar mecanismos e relações capazes de produzir determinados comportamentos, enquanto modelos preditivos têm como objetivo estimar o comportamento futuro ou desconhecido de um sistema sob determinadas condições.

A reprodução de um padrão observado por uma simulação não demonstra, isoladamente, que o modelo possui capacidade preditiva. Da mesma forma, um modelo utilizado para investigar relações entre mecanismos não precisa necessariamente fornecer previsões epidemiológicas para cumprir seu objetivo.

Essa distinção é relevante para a interpretação dos resultados de modelos epidemiológicos. A adequação de um simulador depende de seu propósito, das hipóteses utilizadas, da qualidade dos dados disponíveis e dos procedimentos empregados para sua avaliação.

\section{Modelagem e simulação baseada em agentes}

A modelagem baseada em agentes, ou \textit{Agent-Based Modeling} (ABM), representa sistemas por meio de entidades capazes de possuir atributos, estados e comportamentos próprios. Os agentes interagem entre si e com o ambiente, e essas interações podem produzir padrões coletivos que emergem da dinâmica do sistema \cite{macal_north_2010}.

Essa abordagem é especialmente adequada para sistemas que apresentam heterogeneidade entre entidades e relações dependentes do contexto local. Na representação da dengue, por exemplo, agentes humanos e mosquitos podem possuir localizações e estados epidemiológicos distintos, enquanto o ambiente pode influenciar suas oportunidades de interação.

\subsection{Entidades, atributos e comportamentos}

Um agente pode ser caracterizado por atributos que descrevem suas propriedades e por variáveis de estado que podem ser modificadas durante a execução da simulação. Seus comportamentos correspondem às regras que determinam como o agente interage com outros agentes e com o ambiente.

A definição dessas entidades e regras constitui uma das principais etapas da modelagem baseada em agentes. O comportamento agregado da simulação resulta das interações estabelecidas no nível individual, podendo produzir padrões temporais ou espaciais que não são definidos diretamente como uma única regra global.

Entretanto, a emergência de um padrão na simulação não representa, por si só, uma comprovação de que o modelo reproduz corretamente o sistema real. A interpretação desses padrões depende das hipóteses adotadas e dos procedimentos utilizados para verificar e avaliar o modelo.

\subsection{Interações e vizinhanças}

Em modelos espaciais, as interações entre agentes podem depender de regras de vizinhança ou proximidade. A definição dessas regras determina quais entidades podem interagir em determinado momento e pode influenciar tanto o comportamento do modelo quanto seu custo computacional.

Uma busca exaustiva entre populações de agentes pode exigir um grande número de comparações. Por esse motivo, a representação espacial e as estratégias utilizadas para identificar agentes próximos constituem aspectos relevantes no desenvolvimento de simulações com grande quantidade de entidades.

\section{Modelos epidemiológicos e máquinas de estados}

Modelos epidemiológicos compartimentais representam indivíduos ou populações por meio de estados relacionados à evolução de uma doença. Entre os modelos mais conhecidos encontram-se o SIR, composto pelos estados suscetível, infectado e recuperado, e o SEIR, que acrescenta um estado de exposição ou incubação entre a suscetibilidade e a condição infectante.

Em uma modelagem baseada em agentes, esses estados podem ser representados individualmente. Dessa forma, cada agente possui uma condição epidemiológica própria e realiza transições conforme regras, eventos, temporizadores ou probabilidades definidos pelo modelo.

\subsection{Representação SEIR e SEI}

Para a população humana, uma estrutura SEIR pode representar os estados:

\begin{itemize}
    \item \textbf{Suscetível (S):} indivíduo que pode adquirir a infecção;
    \item \textbf{Exposto (E):} indivíduo infectado que ainda não se encontra em condição infectante;
    \item \textbf{Infectante (I):} indivíduo capaz de participar da transmissão;
    \item \textbf{Recuperado (R):} indivíduo que concluiu o processo epidemiológico representado pelo modelo.
\end{itemize}

Para os mosquitos, uma estrutura SEI pode representar os estados suscetível, exposto e infectante. A utilização desses estados permite distinguir a aquisição da infecção do momento em que o vetor passa a participar efetivamente da transmissão.

\subsection{Máquinas de estados}

A representação computacional de um modelo epidemiológico pode ser compreendida como uma máquina de estados. Nesse caso, é necessário definir quais estados são válidos, quais transições podem ocorrer e quais condições determinam a passagem de um estado para outro.

A especificação dessas transições contribui para a clareza do modelo e para sua posterior verificação. Uma representação conceitual bem definida permite comparar o comportamento esperado com as regras efetivamente implementadas no software.

Além disso, atributos como vacinação ou outras condições individuais podem ser representados separadamente dos estados epidemiológicos. Essa separação permite definir de maneira explícita como tais características influenciam os eventos de infecção ou outros processos representados pela simulação.

\section{Representação espacial e Sistemas de Informação Geográfica}

Os Sistemas de Informação Geográfica permitem organizar e analisar dados associados a localizações no território. Dados vetoriais podem representar pontos, linhas e polígonos, enquanto dados matriciais podem representar superfícies ou variáveis distribuídas espacialmente.

Em uma simulação espacial, os dados geográficos podem desempenhar funções que ultrapassam a visualização do território. As geometrias podem delimitar o ambiente, orientar a distribuição dos agentes, restringir deslocamentos e apoiar cálculos de distância e proximidade.

A integração entre dados geográficos e modelos computacionais exige atenção à compatibilidade entre sistemas de referência de coordenadas, unidades de distância, qualidade das geometrias e níveis de agregação espacial utilizados.

\subsection{Mobilidade e comportamento espacial}

A mobilidade constitui um dos elementos que podem influenciar as oportunidades de interação entre agentes. A escolha de um destino e a definição da trajetória até esse destino representam problemas distintos e podem ser modeladas de diferentes maneiras.

Quando uma simulação utiliza dados agregados, é necessário explicitar quais simplificações são adotadas para representar o movimento individual. Dados sobre deslocamentos entre regiões, por exemplo, não fornecem automaticamente as trajetórias realizadas por cada agente.

Dessa forma, a mobilidade em modelos baseados em agentes deve ser compreendida como uma representação definida a partir dos objetivos do estudo, dos dados disponíveis e das hipóteses adotadas.

\subsection{Integração entre SIG e ABM}

A integração entre Sistemas de Informação Geográfica e modelagem baseada em agentes permite incorporar informações territoriais diretamente ao funcionamento da simulação. Taillandier e Drogoul (2010) discutem diferentes formas de utilização de dados geográficos, incluindo a leitura de arquivos espaciais, o uso das geometrias como elementos que condicionam operações do modelo e a representação de objetos geográficos como agentes.

Na leitura de arquivos espaciais, as geometrias e seus atributos podem ser utilizadas para inicializar ou delimitar o ambiente. Em outras situações, elementos territoriais podem participar diretamente das regras de interação e comportamento. A escolha entre essas possibilidades depende da forma como o espaço é incorporado ao modelo.

Essa integração é particularmente relevante para estudos epidemiológicos que consideram heterogeneidade territorial, distribuição espacial de populações e relações entre localização e interação.

\section{Estocasticidade e experimentação computacional}

Modelos estocásticos incorporam componentes aleatórios em determinados processos. Em uma simulação epidemiológica, a aleatoriedade pode estar presente, por exemplo, na ocorrência de eventos de transmissão, na movimentação dos agentes ou em processos relacionados ao ciclo de vida dos vetores.

A presença de estocasticidade implica que execuções realizadas com os mesmos parâmetros podem produzir resultados diferentes. Por esse motivo, a análise de modelos estocásticos requer procedimentos que considerem a variabilidade entre execuções.

\subsection{Sementes e reprodutibilidade}

Os mecanismos de geração pseudoaleatória utilizam sementes para controlar a sequência de valores produzidos durante uma execução. A utilização de uma mesma configuração e de uma mesma semente permite, em determinadas condições, reproduzir uma sequência de eventos aleatórios.

O controle das sementes contribui para a reprodutibilidade dos experimentos e para a comparação entre diferentes configurações. Entretanto, a reprodução de uma execução também depende da preservação das condições de inicialização, dos parâmetros utilizados e da versão da implementação.

\subsection{Replicações e análise de variabilidade}

Uma única execução de um modelo estocástico pode não representar adequadamente a variabilidade associada aos processos aleatórios. A realização de múltiplas replicações permite observar diferenças entre execuções e calcular medidas descritivas, como médias e dispersões.

A quantidade de replicações e os indicadores utilizados devem ser definidos de acordo com os objetivos do experimento e com o custo computacional das simulações. A comparação entre cenários também deve considerar se as diferenças observadas são consistentes diante da variabilidade entre execuções.

\subsection{Sensibilidade e incerteza}

A análise de sensibilidade investiga como alterações nos parâmetros de um modelo podem influenciar suas saídas. Essa análise pode ser realizada por meio de diferentes estratégias, desde a variação controlada de parâmetros até métodos globais capazes de considerar interações entre múltiplas variáveis.

Os índices de Sobol constituem uma abordagem para analisar a contribuição dos parâmetros para a variância das saídas de um modelo \cite{sobol_2001}. Entretanto, sua aplicação requer um planejamento experimental compatível com o método e um número adequado de execuções.

Em modelos espacialmente explícitos, a incerteza também pode apresentar distribuição heterogênea no território e variar ao longo do tempo. Gamal et al. (2025) discutem uma abordagem para analisar a propagação espaço-temporal da incerteza em modelos baseados em agentes, destacando a importância de considerar tanto a dimensão espacial quanto a evolução temporal dos resultados.

Essas abordagens demonstram que a análise de um modelo estocástico não deve se limitar à observação de uma única execução. A escolha dos procedimentos de experimentação deve, contudo, ser compatível com os objetivos e com os recursos disponíveis para o estudo.

\section{Plataforma GAMA e linguagem GAML}

A plataforma GAMA foi desenvolvida para a construção e execução de modelos espaciais baseados em agentes. A plataforma oferece recursos para a representação de agentes, ambientes geográficos e experimentos computacionais, sendo especialmente voltada a modelos que utilizam informações espaciais \cite{taillandier_etal_2019}.

A linguagem GAML é utilizada para especificar os componentes da simulação. Entre os elementos presentes na linguagem estão estruturas destinadas à definição de variáveis globais, espécies de agentes, procedimentos de inicialização, comportamentos executados durante a simulação e experimentos.

A utilização de uma plataforma especializada facilita a integração entre agentes e informações geográficas, mas não substitui a necessidade de especificar adequadamente as regras do modelo. A qualidade da representação depende das hipóteses adotadas, da implementação e dos procedimentos utilizados para verificar o funcionamento do simulador.

\subsection{Operações espaciais}

Modelos espaciais podem utilizar operações relacionadas ao cálculo de distâncias, à identificação de vizinhos, à manipulação de geometrias e à definição de trajetórias. Esses recursos são relevantes para simulações em que a localização dos agentes influencia suas possibilidades de interação.

Entretanto, a disponibilidade de determinado recurso em uma plataforma não significa que ele seja automaticamente utilizado por todos os modelos desenvolvidos nesse ambiente. A descrição de um simulador específico deve distinguir os recursos disponíveis na plataforma das operações efetivamente incorporadas à sua implementação.

\section{Documentação, verificação e validação}

A construção de um modelo computacional envolve não apenas sua implementação, mas também a documentação das hipóteses e regras utilizadas e a avaliação de seu comportamento.

O protocolo ODD, cujo nome deriva de \textit{Overview, Design concepts and Details}, fornece uma estrutura para documentar modelos baseados em agentes. O protocolo organiza a descrição de aspectos como propósito do modelo, entidades, escalas, processos, condições de inicialização, entradas e submodelos \cite{grimm_etal_2020}.

A utilização de uma estrutura de documentação contribui para tornar explícitas as decisões de modelagem e facilitar a compreensão e a reprodução do trabalho. Entretanto, a documentação, por si só, não comprova que a implementação executa corretamente as regras especificadas.

\subsection{Verificação}

A verificação está relacionada à análise da correspondência entre a implementação e as regras ou especificações definidas para o modelo. Esse processo pode envolver testes de transições de estado, condições de inicialização, regras de interação e indicadores produzidos durante a execução.

A verificação é importante porque uma simulação pode produzir resultados aparentemente plausíveis mesmo quando determinadas regras são implementadas de forma diferente da especificação conceitual.

\subsection{Validação}

A validação está relacionada à adequação do modelo ao propósito para o qual foi desenvolvido e à sua capacidade de representar aspectos relevantes do sistema estudado. Dependendo dos objetivos da pesquisa, esse processo pode envolver a comparação entre resultados simulados e dados observados.

A validação deve ser distinguida da verificação. Um software pode executar corretamente as regras especificadas e, ainda assim, utilizar hipóteses insuficientes para representar adequadamente o fenômeno real. Da mesma forma, uma aproximação visual entre uma saída simulada e uma série observada não comprova isoladamente a correção dos mecanismos internos do modelo.

A calibração também representa um procedimento distinto, relacionado ao ajuste de parâmetros com base em informações observadas.

\subsection{Rastreabilidade e reprodutibilidade}

A rastreabilidade busca estabelecer relações claras entre os elementos conceituais do modelo e sua implementação. Essa relação permite identificar como entidades, estados e regras especificadas são transformados em componentes computacionais.

A reprodutibilidade, por sua vez, depende da documentação das condições de execução, dos parâmetros, das sementes utilizadas, das versões do software e dos procedimentos experimentais. Esses aspectos são particularmente relevantes em modelos estocásticos, nos quais alterações nas condições de execução podem produzir resultados diferentes.

\section{Trabalhos relacionados}

A análise de trabalhos relacionados permite identificar diferentes estratégias utilizadas na aplicação de modelos baseados em agentes a problemas epidemiológicos e espaciais. Os estudos discutidos nesta seção são utilizados como referências para compreender decisões de modelagem, representação espacial e experimentação, sem que seus parâmetros ou resultados sejam automaticamente transferidos para o presente estudo.

\subsection{Garcias (2024)}

Garcias (2024) apresenta um modelo baseado em agentes para representar a transmissão da dengue no contexto de Portugal Continental, utilizando a plataforma GAMA. O estudo considera agentes humanos, mosquitos e componentes territoriais, além de discutir parâmetros, mobilidade e condições ambientais.

Sua principal contribuição para este trabalho está na organização da documentação do modelo e na relação entre representação geográfica, regras de comportamento e experimentação. Entretanto, as diferenças entre a espécie vetorial, a área de estudo e as condições epidemiológicas impedem a transferência automática de parâmetros e resultados.

\subsection{Oliveira et al. (2025)}

Oliveira et al. (2025) apresentam uma simulação multiagentes aplicada à transmissão da dengue no município de Alto Santo, no Ceará, também utilizando a plataforma GAMA. O estudo representa humanos, mosquitos, ovos e criadouros, além de utilizar informações espaciais e realizar experimentos com diferentes configurações de parâmetros.

O trabalho fornece uma referência relevante para a organização de entidades, parâmetros e cenários experimentais. Entretanto, sua área de estudo, suas hipóteses e suas configurações não correspondem diretamente ao contexto do município de São Paulo.

\subsection{Estudos espaciais aplicados a outras doenças}

Gharakhanlou, Mesgari e Hooshangi (2019) e Gharakhanlou, Hooshangi e Helbich (2020) apresentam modelos espaciais baseados em agentes aplicados à transmissão da malária. Embora o sistema epidemiológico seja diferente, esses trabalhos contribuem para a compreensão da integração entre entidades, espaço, fatores ambientais e intervenções.

A principal contribuição desses estudos para o presente trabalho está nos aspectos metodológicos relacionados à representação espacial. As regras epidemiológicas e os parâmetros específicos da malária, entretanto, não podem ser transferidos diretamente para a representação da dengue.

\subsection{Gamal et al. (2025)}

Gamal et al. (2025) apresentam uma abordagem para analisar a propagação espaço-temporal da incerteza em modelos baseados em agentes. O estudo destaca que a variabilidade dos resultados pode apresentar diferenças entre unidades espaciais e ao longo do tempo.

A contribuição do trabalho para esta pesquisa está na discussão sobre a análise de incerteza e heterogeneidade espacial. A abordagem apresentada não constitui um conjunto de regras epidemiológicas para a dengue, mas fornece referências conceituais para a avaliação de modelos espaciais e estocásticos.

\subsection{Síntese dos trabalhos relacionados}

Os trabalhos analisados demonstram diferentes possibilidades de aplicação da modelagem baseada em agentes em contextos epidemiológicos e espaciais. Os estudos relacionados contribuem para a compreensão da documentação de modelos, da representação de entidades, da integração de dados geográficos e da realização de experimentos.

Entretanto, cada modelo está associado a uma combinação específica de objetivos, hipóteses, dados, área de estudo e procedimentos de avaliação. Dessa forma, os resultados e parâmetros apresentados em estudos anteriores não substituem a necessidade de especificação e avaliação próprias do simulador desenvolvido neste trabalho.
